\documentclass{article}

\usepackage{amsmath}
\usepackage{amssymb}
\usepackage{xcolor}
\usepackage{tikz}

\usepackage[utf8]{inputenc}
\usepackage[T1]{fontenc}

% \usepackage[polish]{babel}

\usepackage{listings}
\lstset{
    language=SQL,
    inputencoding=utf8x,
    extendedchars=\true,
    literate={ą}{{\k{a}}}1
             {Ą}{{\k{A}}}1
             {ę}{{\k{e}}}1
             {Ę}{{\k{E}}}1
             {ó}{{\'o}}1
             {Ó}{{\'O}}1
             {ś}{{\'s}}1
             {Ś}{{\'S}}1
             {ł}{{\l{}}}1
             {Ł}{{\L{}}}1
             {ż}{{\.z}}1
             {Ż}{{\.Z}}1
             {ź}{{\'z}}1
             {Ź}{{\'Z}}1
             {ć}{{\'c}}1
             {Ć}{{\'C}}1
             {ń}{{\'n}}1
             {Ń}{{\'N}}1
}

\title{Postać kanoniczna Jordana - przykład}
%\author{You}

\begin{document}
\maketitle

Znaleźć postać kanoniczną Jordana macierzy $A$:
$$
    A
    =
    \begin{bmatrix}\begin{array}{rrrr}
        2   &   0   &   0   &   0\\
        0   &   1   &   1   &   0\\
        0   &   0   &   1   &   0\\
        0   &   2   &  -1   &   2
    \end{array}\end{bmatrix}
$$

Wyznaczamy wielomian charakterystyczny:
$$
    W(\lambda) = |\lambda I - A| = (\lambda - 2)^{2}(\lambda - 1)^{2}
$$

Szukamy wielomianu minimalnego, można to zrobić dzieląc wielomian charakterystyczny, dzieląc go przez największy wspólny dzielnik macierzy dołączonej $adj(\lambda I - A)$:
$$
    adj(\lambda I - A)
    =
    \begin{bmatrix}\begin{array}{cccc}
        (\lambda - 1)^{2}(\lambda - 2)  &   0                               &   0                               &   0\\
        0                               &   (\lambda - 2)^{2}(\lambda - 1)  &   (\lambda - 2)^{2}               &   0\\
        0                               &   0                               &   (\lambda - 2)^{2}(\lambda - 1)  &   0\\
        0                               &   2(\lambda - 1)(\lambda - 2)     &  -1(\lambda - 2)(\lambda - 3)     &   (\lambda - 1)^{2}(\lambda - 2)
    \end{array}\end{bmatrix}
$$
Łatwo zauważyć, że $NWD(adj(\lambda I - A)) = \lambda - 2$.

Zatem wielomian minimalny ma postać:
$$
    M(\lambda) = (\lambda - 2)(\lambda - 1)^{2}
$$
Czyli dzielniki elementarne to: $(\lambda - 2)$, $(\lambda - 2)$, $(\lambda - 1)^{2}$.

Z czego wynika, iż postać kanoniczna Jordana złożona jest z 3 klatek Jordana: 2 - wymiaru $1 \times 1$ związane z dzielnikiem elementarnym $(\lambda - 2)$ - $\begin{bmatrix}\begin{array}{c} 2 \end{array}\end{bmatrix}$, oraz 1 związaną z dzielnikiem elementarnym $(\lambda - 1)^{2}$ wymiaru $1 \times 1$ - $\begin{bmatrix}\begin{array}{cc} 1 & 1\\ 0 & 1 \end{array}\end{bmatrix}$.

Zatem postać kanoniczna Jordana ma ostatecznie postać:
$$
    J
    =
    \begin{bmatrix}\begin{array}{rrrr}
        1   &   1   &   0   &   0\\
        0   &   1   &   0   &   0\\
        0   &   0   &   2   &   0\\
        0   &   0   &   0   &   2
    \end{array}\end{bmatrix}
$$

Dodatkowo można wyznaczyć macierz modalną:
\begin{itemize}
    \item Wektory związane z $\lambda_{1} = 1$:
    \begin{itemize}
        \item
        $$
        \begin{array}{rrrr|r}
           -1   &   0   &   0   &   0   &   0\\
            0   &   0   &  -1   &   0   &   0\\
            0   &   0   &   0   &   0   &   0\\
            0   &  -2   &   1   &  -1   &   0
        \end{array}
        \rightarrow
        h_{1}^{1}
        =
        \begin{bmatrix}\begin{array}{c}
            0\\
            a\\
            0\\
            -2a
        \end{array}\end{bmatrix}
        \rightarrow
        h_{1}^{1}
        =
        \begin{bmatrix}\begin{array}{c}
            0\\
            1\\
            0\\
            -2
        \end{array}\end{bmatrix}
        $$
        \item
        $$
        \begin{array}{rrrr|c}
           -1   &   0   &   0   &   0   &   0\\
            0   &   0   &  -1   &   0   &   a\\
            0   &   0   &   0   &   0   &   0\\
            0   &  -2   &   1   &  -1   &   -2a
        \end{array}
        \rightarrow
        h_{1}^{2}
        =
        \begin{bmatrix}\begin{array}{c}
            0\\
            a\\
            b\\
            -b -2a
        \end{array}\end{bmatrix}
        \rightarrow
        h_{1}^{2}
        =
        \begin{bmatrix}\begin{array}{c}
            0\\
            0\\
            1\\
            -1
        \end{array}\end{bmatrix}
        $$
    \end{itemize}

    \item Wektory związane z $\lambda_{2} = 2$:
        $$
        \begin{array}{rrrr|r}
            0   &   0   &   0   &   0   &   0\\
            0   &   1   &  -1   &   0   &   0\\
            0   &   0   &   1   &   0   &   0\\
            0   &  -2   &   1   &   2   &   0
        \end{array}
        \rightarrow
        h_{2}
        =
        \begin{bmatrix}\begin{array}{c}
            b\\
            0\\
            0\\
            c
        \end{array}\end{bmatrix}
        \rightarrow
        h_{2}^{1}
        =
        \begin{bmatrix}\begin{array}{c}
            1\\
            0\\
            0\\
            0
        \end{array}\end{bmatrix}
        ,\,
        h_{2}^{2}
        =
        \begin{bmatrix}\begin{array}{c}
            0\\
            0\\
            0\\
            1
        \end{array}\end{bmatrix}
        $$
\end{itemize}
\end{document}